\chapter{Entrevista con experto}

La guía solicita introducir en ChatGPT un prompt específico para conversar con Carlos Pérez, trabajador con más de diez años de experiencia en Alestis Puerto Real. Se presenta una síntesis de la conversación usada para extraer requisitos: tono informal, alguna duda razonable, referencias puntuales a Cádiz y foco principal en el proceso productivo. Las capturas reales de la conversación deben incorporarse en el Anexo A antes de subir la entrega al campus.

\begin{interviewbrief}
\textbf{Ficha de la entrevista}

\begin{tabularx}{\textwidth}{@{}P{0.28\textwidth}Y@{}}
\textbf{Entrevistado} & Carlos Pérez, operario/técnico con experiencia en operaciones de Alestis Puerto Real.\\
\textbf{Rol en la actividad} & Experto de planta consultado por Lunabotics para identificar problemas internos y oportunidades de automatización.\\
\textbf{Temas tratados} & HTP del A320, sección 19.1, flujo de kits y herramientas, FOD, seguimiento de materiales, transporte interno, esperas y restricciones de seguridad.\\
\textbf{Uso de la evidencia} & Las respuestas se usan para definir requisitos, selección de robot, sensores, arquitectura, presupuesto y KPIs.\\
\textbf{Capturas} & El Anexo A carga automáticamente las imágenes \texttt{captura\_01.png} a \texttt{captura\_04.png} cuando se guarden en \path{figures/entrevista/}.\\
\end{tabularx}
\end{interviewbrief}

\section{Transcripción de la entrevista}

\begin{interviewtranscript}[Transcripción de trabajo con Carlos Pérez]
\InterviewExchange
{Carlos, para situarnos: ¿cómo explicarías el HTP y la sección 19.1 sin ponerte demasiado de libro?}
{El HTP es la cola del avión, el estabilizador horizontal. En planta lo llamamos así, HTP, y no es una pieza pequeña precisamente. La sección 19.1 va por la parte posterior, cerca de las interfaces de cola. Allí hay muchas operaciones encadenadas. Si una se retrasa, la siguiente lo nota enseguida. Desde fuera parece que todo es montar piezas, pero alrededor de la estructura se mueve media nave.}

\InterviewExchange
{Cuando dices que se mueve media nave, ¿qué cosas se mueven de verdad durante un turno?}
{Uf, muchísimo más de lo que parece. Carros con utillaje, piezas pendientes de montar, material de sellado, consumibles... y cada carro con su identificación y su sitio, porque si no aquello sería un caos. También se mueven herramientas controladas: taladros, calibres, dinamométricas. Hay días en los que parece que media nave busca una herramienta ``desaparecida'' y luego aparece debajo de una funda o dentro de un útil. Y eso enlaza con FOD, claro. Cualquier cosa fuera de control obliga a revisar.}

\InterviewExchange
{Además de carros y herramientas, ¿hay mucho movimiento de documentación o de personas?}
{Sí, aunque ahora se ve más en tablets y terminales. Siguen estando las órdenes de trabajo, no conformidades, firmas de calidad, validaciones de inspección... y alguna carpeta vieja todavía ronda por ahí. Luego está la gente: calidad entrando y saliendo, logística trayendo material, producción preguntando por avances, mantenimiento revisando útiles. En la 19.1 eso pesa bastante. Hay turnos en los que el verdadero trabajo es coordinar el flujo alrededor de la pieza, más que tocar la pieza en sí.}

\InterviewExchange
{¿Dónde se va más tiempo: en montaje, en esperas o en retrabajos?}
{Depende del día y del programa, pero si preguntas dentro te dirán casi lo mismo: esperas y retrabajos. Esperar una inspección, una firma de calidad, una pieza que ``ya venía de camino''... eso mata el ritmo. Y en aeroestructuras no puedes decir ``sigo y luego lo vemos''. Si falta una validación, se para. Los retrabajos también duelen: un agujero fuera de tolerancia, un remache mal puesto, sellante donde no debía. Desmontar, limpiar, volver a hacer. Ahí se van horas.}

\InterviewExchange
{¿Y el FOD? A veces se nombra como si fuera solo orden y limpieza.}
{FOD es serio. Una punta de broca, un tornillo, una brida, un trozo de embalaje. Cuando falta una herramienta o no cuadra un consumible, puedes acabar parando una zona hasta asegurarte de que no se quedó nada dentro de la estructura. He visto búsquedas larguísimas por cosas diminutas. No siempre pasa, pero cuando pasa consume tiempo y pone nervioso a todo el mundo.}

\InterviewExchange
{Entonces, ¿tendría sentido que un robot móvil fabricara o montara directamente el HTP?}
{Fabricarlo entero, yo diría que no. Al menos ahora mismo. En el montaje hay procedimientos muy estrictos, gente cualificada y mucho ajuste fino. A veces el operario nota a tacto si una pieza asienta bien, si un borde está raro o si el sellante no quedó como debía. Un robot necesitaría una barbaridad de sensores y aun así tendría limitaciones. Donde sí lo veo es alrededor: mover carros, llevar herramientas o consumibles, hacer una ronda visual básica, escanear FOD, registrar entregas. Pero montaje completo del HTP, complicado.}

\InterviewExchange
{¿El entorno de Cádiz influye en cómo se trabaja y se convive en planta?}
{Sí, se nota. Hay mucha cultura de oficio en la Bahía, mucha gente que viene de astilleros, metal, aeronáutica o de familias que llevan años en esto. El trato suele ser cercano, con bastante compañerismo, aunque cuando toca auditoría o FOD la gente se pone seria. También pasa que viene gente de visita, ve una industria muy avanzada y se sorprende de cuánto trabajo manual sigue habiendo. Luego salen a Cádiz, se comen un pescaito y dicen que lo mejor fue la Bahía. Tampoco les culpo. Para el robot eso importa: si vendes que va a hacerlo todo, nadie te cree; si dices que va a quitar paseos, ordenar entregas y dejar registro, suena más real.}

\InterviewExchange
{Si el robot se queda en tareas de soporte, ¿qué requisitos mínimos tendría que cumplir?}
{Bastantes más de lo que parece. No es como meterlo en un almacén con pasillos amplios y cajas iguales. En la 19.1 hay útiles enormes, plataformas, gente trabajando alrededor y material temporal. El AMR tendría que ir despacio, detectar personas agachadas, herramientas fuera de sitio y obstáculos raros. También debería cuidar no generar FOD: ruedas, piezas sueltas, vibraciones, mantenimiento. Y, por supuesto, registrar qué recogió, dónde lo dejó, a qué hora y quién confirmó la entrega.}

\InterviewExchange
{¿Un AGV de ruta fija bastaría?}
{Depende de lo que quieras resolver. Para logística muy perimetral o rutas controladas puede servir. Pero en montaje se queda corto rápido: cambian carros, útiles, plataformas, inspecciones que bloquean accesos. Una ruta fija se vuelve rígida. Yo preferiría un robot con mapa, zonas restringidas y capacidad de buscar alternativa. Si el pasillo está bloqueado de verdad, esperará o avisará; milagros no hace.}

\InterviewExchange
{¿Cree que la aceptación de operarios depende de cómo se presente el proyecto?}
{Muchísimo. Si se presenta como ``esto viene a sustituir tareas'' o como si la gente trabajara lento, la reacción será defensiva. Si se plantea como apoyo, menos paseos, menos búsqueda de carros y menos estrés logístico, cambia bastante. También importa quién lo presenta. No es lo mismo alguien de producción que pisa planta todos los días que un equipo externo que llega con un PowerPoint. Si el operario ve que le facilita el turno, lo adopta; si añade pasos o burocracia, lo cumple lo mínimo.}

\InterviewExchange
{La actividad habla de pasar de unos 40 HTP al mes a 80 o más. ¿Cómo se debe leer esa meta?}
{No diría que el robot duplica producción. Eso sería pasarse. Ayuda en cosas concretas: menos esperas, menos búsquedas, entregas más fiables y mejor registro. Para llegar a 80 hacen falta planificación, gente, calidad y capacidad de estación. El robot aporta una parte, no toda la solución.}

\InterviewExchange
{Si Lunabotics empieza con un piloto, ¿qué mediría antes de comprar 7 o 17 robots?}
{Misiones completadas, entregas a tiempo, paradas de seguridad, lecturas QR/RFID correctas, disponibilidad y eventos FOD revisados. Y algo muy simple: si el operario camina menos de verdad. Si eso no sale, yo no compraría siete robots. Primero que el piloto demuestre que sirve.}
\end{interviewtranscript}

\section{Hallazgos para diseño}

La entrevista apunta a varias decisiones de diseño. El robot debe trabajar en logística auxiliar: carros, kits, útiles, herramientas y consumibles. El registro tiene que ocurrir en origen y destino, porque saber que el material existe no alcanza para saber que está en la estación. FOD debe tratarse como inspección preventiva con revisión humana. Y el piloto debe medirse con datos de operación antes de hablar de flotas grandes.

\begin{longtable}{P{0.42\textwidth}P{0.50\textwidth}}
\toprule
\textbf{Lo que dice planta} & \textbf{Decisión de diseño}\\
\midrule
\endhead
Carros, herramientas y consumibles se mueven durante todo el turno. & AMR para entregas, retornos y reposición de kits/herramientas.\\
Herramientas y útiles pueden no estar donde el sistema cree. & RFID/QR en recogida, entrega y devolución, ligado a misión y estación.\\
Esperas, firmas y retrabajos rompen el ritmo de la sección. & Misiones programadas por ventana horaria, confirmación HMI y registro de incidencias.\\
El robot será aceptado si no estorba y si resuelve trabajo real. & Piloto pequeño, rutas validadas, HMI sencilla y medición antes/después.\\
FOD exige evidencia, pero la decisión sigue siendo de calidad. & Cámara RGB-D/industrial para ronda preventiva y registro, con revisión humana.\\
La planta cambia durante el turno. & AMR omnidireccional con SLAM, replanificación y gestor de misiones para evitar rutas fijas.\\
\bottomrule
\end{longtable}

\section{Lectura para Lunabotics}

La entrevista deja claro que la automatización tiene sentido en suministro, retorno, registro y rondas preventivas. En nuestra evaluación, el robot debe demostrar valor en una zona reducida antes de convertirse en flota. Ese enfoque baja el riesgo comercial porque obliga a medir antes de ampliar.
